% status: experimental --- NOT part of the published results
% Input-ready body; requires only amsmath, amssymb, and amsthm.
% M14 is reserved for integration by the coordinating author.
\section{M14: a controlled remainder for the gauge two-term law}
\label{sec:two-term-remainder}

\subsection{Setting and an exact nonnegative remainder}
\label{subsec:m14-identity}

All matrices below refer to the fixed, whitened linearization, not to a
nonlinear reconstruction ensemble.  Let
$A\in\mathbb{R}^{n\times p}$ have full column rank,
$B\in\mathbb{R}^{n\times m}$, and $\Sigma_0\succ0$.  Write
\begin{equation}
 D=A^\top A,\quad U=A^\top B,\quad M=B^\top B,\quad
 \Lambda_0=\Sigma_0^{-1},\qquad
 F(t)=D-U(M+t\Lambda_0)^{-1}U^\top,\quad t>0.
 \label{eq:m14-fisher}
\end{equation}
The precision multiplier $t$ is uniform on the nuisance coordinates included
in this model; inactive lights with zero Jacobians can be omitted.  Fix a
unit Euclidean task vector $a\in\mathbb{R}^p$ and a specified gauge vector
$c\in\mathbb{R}^m$ satisfying
\begin{equation}
 \|a\|=1,\qquad Aa=Bc.
 \label{eq:m14-gauge}
\end{equation}
In the fixed-normal albedo experiment, $a=\rho/\|\rho\|$ and
$c_k=(1,0,0)^\top/\|\rho\|$ for each active light.  These are the coordinates
of the actual Jacobians in \eqref{eq:m14-fisher}; no change to a log-albedo
coordinate system is made here.  Define
\begin{equation}
 J_a(t)=a^\top F(t)^{-1}a,\qquad
 q=c^\top\Lambda_0c>0,\qquad r=\frac{1}{\|Aa\|^2},\qquad
 j_t=\frac{1}{tq}+r.
 \label{eq:m14-main}
\end{equation}
The choice of $c$ matters if several nuisance vectors realize the same gauge;
$c$ is not silently replaced by a least-squares or minimum-prior-norm gauge.

Here are all ingredients of the remainder.  Put
\begin{equation}
 P_A=AD^{-1}A^\top,\quad T=B^\top(I-P_A)B=M-U^\top D^{-1}U,\quad
 g=U^\top D^{-1}a.
 \label{eq:m14-projection}
\end{equation}
Choose any nonsingular factor $C$ with $\Sigma_0=CC^\top$ (a Cholesky factor
is sufficient), and set
\begin{align}
 S&=C^\top T C, & h&=C^{-1}c/\sqrt q, & v&=C^\top g,\\
 e_0&=v-h/\sqrt q=C^\top g-C^{-1}c/q, &
 Q^\top Q&=I_{m-1}, & QQ^\top&=I-hh^\top,\\
 S_\perp&=Q^\top S Q, & e&=Q^\top e_0, &
 \delta_D&=a^\top D^{-1}a-r.
 \label{eq:m14-residual-definitions}
\end{align}
For $m=1$ the perpendicular matrices and vectors are empty, and their
quadratic forms are zero.  Other null vectors of $T$, if present, are
\emph{retained} in $S_\perp$; only the named gauge $h$ is removed.

\paragraph{Proposition M14.1 (exact decomposition).}
Under \eqref{eq:m14-fisher}--\eqref{eq:m14-gauge}, $F(t)\succ0$ and
\begin{equation}
 \boxed{\quad J_a(t)=j_t+\mathcal R_t,\qquad
 \mathcal R_t=\delta_D+e^\top(tI+S_\perp)^{-1}e\geq0.\quad}
 \label{eq:m14-exact-remainder}
\end{equation}
In particular, an exact gauge gives a one-sided approximation, not a
small-error theorem.  Equality $J_a(t)=j_t$ at any finite $t>0$ holds if and
only if both $Da=(a^\top Da)a$ and $g=\Lambda_0c/q$.

\begin{proof}
The joint quadratic form
$\|Ax+By\|^2+t y^\top\Lambda_0y$ is positive for every nonzero pair
$(x,y)$: if it were zero, positive definiteness of $\Lambda_0$ gives $y=0$,
and the full column rank of $A$ then gives $x=0$.  Its two Schur complements
are therefore positive definite.  Block elimination gives
\begin{equation}
 F(t)^{-1}=D^{-1}+D^{-1}U(t\Lambda_0+T)^{-1}U^\top D^{-1}.
 \label{eq:m14-pushthrough}
\end{equation}
$P_A$ is an orthogonal projector, so $T\succeq0$.  The gauge implies
$Tc=B^\top(I-P_A)Aa=0$ and
$c^\top g=(Aa)^\top AD^{-1}a=a^\top a=1$.
Consequently, $\|h\|=1$, $Sh=0$, and $h^\top v=1/\sqrt q$.
Since $S$ is symmetric, $h^\perp$ is invariant under $S$, and
\begin{equation}
 (t\Lambda_0+T)^{-1}=C(tI+S)^{-1}C^\top,\qquad
 v^\top(tI+S)^{-1}v=\frac{1}{tq}
                   +e^\top(tI+S_\perp)^{-1}e.
\end{equation}
Contracting \eqref{eq:m14-pushthrough} with $a$ proves the identity.
Cauchy--Schwarz applied to $D^{1/2}a$ and $D^{-1/2}a$ gives
$(a^\top Da)(a^\top D^{-1}a)\geq1$, hence $\delta_D\geq0$.
The other remainder term is nonnegative because $tI+S_\perp\succ0$.
The equality condition for Cauchy--Schwarz is $Da=(a^\top Da)a$.
The positive definite quadratic form vanishes exactly when $e=0$.
As $e_0\perp h$, this is equivalent to $e_0=0$, or
$C^\top g=C^{-1}c/q$, equivalently $g=\Lambda_0c/q$.
\end{proof}

The residual norm and the diagonal part have useful measurable forms:
\begin{align}
 \|e\|^2
 &=g^\top\Sigma_0g-\frac1q
   =\left(g-\frac{\Lambda_0c}{q}\right)^\top
       \Sigma_0\left(g-\frac{\Lambda_0c}{q}\right),
 \label{eq:m14-prior-residual}\\
 \delta_D
 &=\frac{a^\top(D-\mu I)D^{-1}(D-\mu I)a}{\mu^2},\qquad
 \mu=a^\top Da.
 \label{eq:m14-diagonal-residual}
\end{align}
These identities follow by expansion using $c^\top g=1$ and $a^\top a=1$.
For diagonal $D=\operatorname{diag}(d_i)$, the last expression is the
cancellation-free sum $\sum_i a_i^2(d_i-\mu)^2/(\mu^2d_i)$.
Thus the two distinct error mechanisms are spatial information variation
and a prior-weighted gauge residual.  Neither is controlled by $Aa=Bc$ alone.

For clarity, the scalar information seen after eliminating $a^\perp$ is
\begin{equation}
 s_a(t)=a^\top F(t)a-f_\times^\top F_\perp(t)^{-1}f_\times
       =\frac1{j_t+\mathcal R_t},\qquad
 \frac1{j_t}-s_a(t)=\frac{\mathcal R_t}{j_t(j_t+\mathcal R_t)}.
 \label{eq:m14-scalar-schur}
\end{equation}
Here $[a,Q_a]$ is orthogonal, $f_\times=Q_a^\top F(t)a$, and
$F_\perp=Q_a^\top F(t)Q_a$.  This follows by the scalar block inverse formula.
It is an effective-information expansion, not a claim that a scalar
approximation controls every entry or every mode of $F(t)$.

\subsection{A non-circular spectral certificate}
\label{subsec:m14-certificate}

Let $S_\perp z_i=\lambda_i z_i$ with orthonormal $z_i$, and
$w_i=(z_i^\top e)^2\geq0$.  Then
$\mathcal R_t=\delta_D+\sum_i w_i/(t+\lambda_i)$.
Instead of evaluating these exact denominators and calling the answer a
bound, partition the nonnegative spectrum into intervals
$[\ell_b,u_b]$ and retain only their masses
$m_b=\sum_{i\in b}w_i$.  The certificate takes
$(q,r,\delta_D,\{\ell_b,u_b,m_b\})$, not $J_a$ or an observed error.

\paragraph{Proposition M14.2 (endpoint interval and factor-two certificate).}
For all $t>0$ define
\begin{equation}
 L_t=\delta_D+\sum_b\frac{m_b}{t+u_b},\qquad
 U_t=\delta_D+\sum_b\frac{m_b}{t+\ell_b}.
 \label{eq:m14-bins}
\end{equation}
Then $L_t\leq\mathcal R_t\leq U_t$.  With the fixed dyadic choice
\begin{equation}
 \ell_b=2^b-1,\quad u_b=2^{b+1}-1,\quad
 b=\lfloor\log_2(1+\lambda_i)\rfloor,
 \label{eq:m14-dyadic}
\end{equation}
we moreover have, uniformly over all admissible models,
\begin{equation}
 \mathcal R_t\leq U_t\leq2\mathcal R_t\qquad (t\geq1).
 \label{eq:m14-factor-two}
\end{equation}
If $\mathcal R_t=0$, both sides are zero and a bound/error ratio is undefined,
not assigned a favorable value.

\begin{proof}
For each eigenvalue in bin $b$,
$1/(t+u_b)\leq1/(t+\lambda_i)\leq1/(t+\ell_b)$.
Multiplication by $w_i\geq0$ and summation gives
\eqref{eq:m14-bins}.  The dyadic bins satisfy
$u_b=2\ell_b+1$.  For $t\geq1$,
\[
 \frac{t+\lambda_i}{t+\ell_b}
 \leq\frac{t+u_b}{t+\ell_b}
 =2+\frac{1-t}{t+\ell_b}\leq2.
\]
Apply this term by term to $U_t$ and use $\delta_D\geq0$.
\end{proof}
For a caller-supplied non-dyadic bin summary, a factor-two assertion instead
requires $(t+u_b)/(t+\ell_b)\leq2$ in every positive-mass bin.  The public
implementation checks these actual ratios; $t\geq1$ alone is not sufficient
for arbitrary intervals such as $[0,100]$.

This is an a posteriori certificate from the nominal design, with an
$m-1$ dimensional spectral calculation.  It does not assert a computational
speed advantage over an exact nuisance-space risk calculation.  Its
non-circularity is substantive: moving eigenvalues inside their bins while
holding the spectral masses fixed changes the true remainder but leaves the
certificate unchanged.  The bin width is fixed by
\eqref{eq:m14-factor-two}, not chosen after observing errors.

The following are directly testable sufficient conditions:
\begin{align}
 U_t\leq\eta j_t&\quad\Longrightarrow\quad
      0\leq\frac{J_a(t)-j_t}{j_t}\leq\eta,\\
 \frac{J_a(t)-j_t}{J_a(t)}&\leq\frac{U_t}{j_t+U_t},\\
 \mathcal R_t&\leq\delta_D+\frac{\|e\|^2}{t+\lambda_*}
 \quad\hbox{if }S_\perp\succeq\lambda_*I,\quad\lambda_*\geq0.
 \label{eq:m14-sufficient}
\end{align}
The first and last statements follow from the endpoint bound and spectral
order; the middle statement follows because $x/(j_t+x)$ increases for
$x\geq0$.  In particular $\lambda_*=0$ always works, although it can be much
looser than the binned bound.  For diagonal $D$, if
$|d_i/\mu-1|\leq\xi<1$ on the support of $a$, then
\begin{equation}
 \delta_D\leq r\frac{\xi^2}{1-\xi},\qquad
 r\frac{\xi^2}{1-\xi}+\frac{\|e\|^2}{t+\lambda_*}\leq\eta j_t
 \quad\Longrightarrow\quad\frac{J_a(t)-j_t}{j_t}\leq\eta.
 \label{eq:m14-physical-sufficient}
\end{equation}
Indeed, substitute $d_i=\mu(1+z_i)$ in
\eqref{eq:m14-diagonal-residual} and bound
$z_i^2/(1+z_i)\leq\xi^2/(1-\xi)$.  Strongly nonuniform scenes need not meet
this coarse condition; the exact positive sum and spectral-bin certificate
remain valid without $\xi<1$.

\subsection{A joint-endpoint certificate for calibration value}
\label{subsec:m14-value}

Let $\kappa>1$, $V=1-J_a(\kappa)/J_a(1)$, and
\begin{equation}
 \beta=\frac{j_\kappa}{j_1}
       =\frac{1/\kappa+qr}{1+qr},\qquad
 V_{\rm pred}=1-\beta=\frac{1-1/\kappa}{1+qr}.
 \label{eq:m14-value-prediction}
\end{equation}
One must not presume that the endpoint errors are independent.  Their exact
coupled correction is
\begin{align}
 V-V_{\rm pred}&=-\frac{N}{j_1+\mathcal R_1},\\
 N&=(1-\beta)\delta_D+\sum_i w_i f(\lambda_i),\qquad
 f(x)=\frac1{\kappa+x}-\frac{\beta}{1+x}
     =\frac{(1-\beta)(x-\kappa qr)}{(1+x)(\kappa+x)}.
 \label{eq:m14-value-correction}
\end{align}
This follows by substituting $J_a(t)=j_t+\mathcal R_t$ and collecting
$\mathcal R_\kappa-\beta\mathcal R_1$.

\paragraph{Proposition M14.3 (signed value interval).}
For each spectral bin let
\begin{align}
 f_b^-&=\min\{f(\ell_b),f(u_b)\},\\
 f_b^+&=\max\bigl(\{f(\ell_b),f(u_b)\}
               \cup\{f(x_*):\ell_b\leq x_*\leq u_b\}\bigr),\qquad
 x_* = \frac{\kappa\sqrt\beta-1}{1-\sqrt\beta},\\
 N_-&=(1-\beta)\delta_D+\sum_b m_b f_b^-,\quad
 N_+=(1-\beta)\delta_D+\sum_b m_b f_b^+,\\
 D_-&=j_1+L_1>0,\qquad D_+=j_1+U_1,\qquad
 \mathcal C=\{n/d:n\in\{N_-,N_+\},\ d\in\{D_-,D_+\}\}.
 \label{eq:m14-value-interval}
\end{align}
Then
\begin{equation}
 -\max\mathcal C\ \leq V-V_{\rm pred}\leq-\min\mathcal C,
 \qquad |V-V_{\rm pred}|\leq
 B_V:=\max_{z\in\mathcal C}|z|.
 \label{eq:m14-value-bound}
\end{equation}
Thus $B_V\leq\varepsilon$ certifies absolute accuracy $\varepsilon$;
if the entire signed interval lies above $\varepsilon$ or below
$-\varepsilon$, it instead certifies failure of that accuracy target.

\begin{proof}
Since $1/\kappa<\beta<1$, the derivative
$f'(x)=\beta/(1+x)^2-1/(\kappa+x)^2$ has exactly one zero on
$[0,\infty)$, namely $x_*>0$.  It is positive before $x_*$ and negative
afterwards, so this is the unique maximum and an interval minimum is at an
endpoint.  It follows that $f_b^-\leq f(\lambda_i)\leq f_b^+$ within bin $b$.
Nonnegative masses give $N_-\leq N\leq N_+$.
Proposition M14.2 gives $D_-\leq j_1+\mathcal R_1\leq D_+$.
For a positive denominator, the extrema of $n/d$ on this rectangle occur
among its four corners (monotonicity in $n$, and the sign-dependent
monotonicity in $d$).  The negative sign in
\eqref{eq:m14-value-correction} yields \eqref{eq:m14-value-bound}.
\end{proof}

There is \emph{no universal relative-tightness assertion for $B_V$}.
For example, take
\[
 A=\begin{pmatrix}1\\0\end{pmatrix},\quad
 B=\begin{pmatrix}1&1\\0&\sqrt{10}\end{pmatrix},\quad
 a=1,\quad c=\begin{pmatrix}1\\0\end{pmatrix},\quad\Sigma_0=I_2.
\]
Then $q=r=1$, $\delta_D=0$, $S_\perp=10$, $e=1$, and
$J_a(t)=1+1/t+1/(t+10)$.  At $\kappa=10$, the nonzero endpoint remainders
cancel exactly in \eqref{eq:m14-value-correction}, giving
$V=V_{\rm pred}=9/20$.  The dyadic bin $[7,15]$ nevertheless gives $B_V>0$.
Perturbing the squared lower-right entry of $B$ from $10$ to $10+u$ makes
$V-V_{\rm pred}=-9u/9200+O(u^2)$, while the bin certificate stays bounded
away from zero.  Thus an arbitrarily large bound/observed ratio for value
is possible, despite the uniform factor-two endpoint result.  The finite
benchmark ratios below are numerical validation, not a proof eliminating
this obstruction.

\subsection{Physical interpretation, heterogeneity, and coupling}
\label{subsec:m14-physical}

For the per-light model let $R_\rho=\|\rho\|$, $L$ be the active-light count,
and
\[
 d_p=\sum_k w_{kp}s_{kp}^2,\qquad
 \ell_k=\sum_p a_p^2\frac{w_{kp}s_{kp}^2}{d_p},\qquad
 g_k=(R_\rho\ell_k,g_{k1},g_{k2})^\top.
\]
The equality for the intensity coordinate follows from
$B_{k0,p}=\rho_p s_{kp}$; $\ell_k\geq0$ and $\sum_k\ell_k=1$.
Here $g_{k1},g_{k2}$ are computed from the physical normal--tangent
Jacobians and weights, so they retain the effects of normals, shadows,
texture, and noise precision.  Suppose first that the channels are uncoupled,
with per-light standard deviations $s_{Ik}>0$ for log intensity and
$s_{Dk}>0$ radians for each tangent coordinate.  Define
\begin{equation}
 \pi_k=\frac{s_{Ik}^{-2}}{\sum_j s_{Ij}^{-2}},\qquad
 q=\frac{\sum_k s_{Ik}^{-2}}{R_\rho^2}.
\end{equation}
Equation \eqref{eq:m14-prior-residual} becomes the exact positive decomposition
\begin{equation}
 \|e\|^2=
 R_\rho^2\sum_k s_{Ik}^2(\ell_k-\pi_k)^2
 +\sum_k s_{Dk}^2(g_{k1}^2+g_{k2}^2).
 \label{eq:m14-uncoupled-physical}
\end{equation}
To prove it, observe that $(\Lambda_0c/q)_{k0}=R_\rho\pi_k$ and its
direction coordinates are zero, then use the block-diagonal covariance
in \eqref{eq:m14-prior-residual}.  In particular, homogeneous channels give
$R_\rho^2s_I^2\sum_k(\ell_k-1/L)^2+
 s_D^2\sum_k(g_{k1}^2+g_{k2}^2)$.
Even with no direction uncertainty, heterogeneous prior precision can
mismatch the illumination leverages $\ell_k$.

The C9 family uses $s_{Ik}=\bar s_I\exp(\tau z_k)$ and
$s_{Dk}=\bar s_D\exp(\tau z_k)$ with a fixed-seed vector $z_k$, and blocks
\begin{equation}
 \Sigma_k=\begin{pmatrix}
 s_{Ik}^2 & \varrho s_{Ik}s_{Dk}/\sqrt2 & \varrho s_{Ik}s_{Dk}/\sqrt2\\
 \varrho s_{Ik}s_{Dk}/\sqrt2 &s_{Dk}^2&0\\
 \varrho s_{Ik}s_{Dk}/\sqrt2 &0&s_{Dk}^2
 \end{pmatrix},\qquad |\varrho|<1.
 \label{eq:m14-coupled-family}
\end{equation}
Here $\tau$ is the heterogeneity parameter and $\varrho$ the correlation
between log intensity and the normalized sum of the tangent coordinates.
In particular,
\begin{equation}
 q=\frac1{R_\rho^2}\sum_k\frac1{s_{Ik}^2(1-\varrho^2)},\qquad
 d_k^g=g_k-\frac{\Sigma_k^{-1}(1,0,0)^\top}{R_\rho q},\qquad
 \|e\|^2=\sum_k(d_k^g)^\top\Sigma_k d_k^g.
 \label{eq:m14-coupled-physical}
\end{equation}
The first formula follows by a scalar Schur complement of the direction
block; the last follows from \eqref{eq:m14-prior-residual}.  The normalized
correlation matrix in \eqref{eq:m14-coupled-family} has eigenvalues
$1$, $1-\varrho$, $1+\varrho$.  Consequently
\[
 \|e\|^2\leq(1+|\varrho|)\sum_k
 \left[s_{Ik}^2(d_{k0}^g)^2+s_{Dk}^2\{(d_{k1}^g)^2+(d_{k2}^g)^2\}\right]
\]
is a purely physical sufficient input to \eqref{eq:m14-sufficient}.

A standard deviation reported as degrees in the C9 input is per tangent
coordinate: $s_{Dk}=(\pi/180)\sigma_{\mathrm{dir},k}^{\mathrm{deg}}$.
The two-coordinate first-order angular RMS is $\sqrt2$ times this per-axis
standard deviation.  It is not identified here with any different radial
rotation convention of a nonlinear corruption sampler.  The artifact
reports both quantities.  Thus \eqref{eq:m14-sufficient} and
\eqref{eq:m14-coupled-physical} are measurable conditions in degrees,
log-intensity deviation, heterogeneity, coupling, and nominal Jacobians.
They do not reduce to a geometry-free law in $s_D/s_I$: the $\delta_D$,
illumination-leverage, and normal--tangent terms would then be missing.

\subsection{Why direction-only refinement can invalidate the two-term law}
\label{subsec:m14-direction}

Keep $s_D>0$ and let the uncoupled intensity deviation $s_I$ tend to zero.
For every $s_I>0$ the theorem applies, $q=L/(R_\rho^2s_I^2)$ diverges,
$j_t\to r$, and $V_{\rm pred}\to0$.  By
\eqref{eq:m14-uncoupled-physical}, the direction part of $\|e\|^2$ need not
tend to zero.  Its contribution to $\mathcal R_t$ can remain strongly
$t$-dependent, giving nonzero $V$.  Exactly at $s_I=0$, the full covariance
is singular and $q$ is not defined by a finite precision inverse.
\emph{Its Moore--Penrose inverse must not be substituted for the precision:}
a known intensity coordinate means infinite precision, not zero precision.
One must either take the positive-definite limit or remove the known
coordinate and rederive the information model.  The reduced direction-only
model need not have the specified intensity gauge at all.

This mechanism occurs within a legal Lambertian block, even at a small
angular standard deviation.  Take one pixel, one light, $\rho=1$,
$n=(4/5,0,3/5)$, $d=(0,0,1)$, tangent vectors $(1,0,0)$ and $(0,1,0)$,
and $w=10000$.  The pixel is unshadowed, with $n^\top d=3/5$, and
\begin{equation}
 A=[60],\quad B=[60,80,0],\quad a=1,\quad c=(1,0,0)^\top,\quad
 \Sigma_0=\operatorname{diag}(\epsilon^2,1/10000,1/10000),\quad\epsilon>0.
 \label{eq:m14-lambertian-counterexample}
\end{equation}
These are exactly $A=\sqrt w\,s$ and
$B=\sqrt w\,\rho(s,n^\top t_1,n^\top t_2)$, not arbitrary Fisher entries.
The gauge is exact; direct scalar inversion gives
\begin{equation}
 J_a(t)=\frac1{3600}+\frac{\epsilon^2}{t}+\frac4{22500t},\qquad
 j_t=\frac1{3600}+\frac{\epsilon^2}{t},\qquad
 \mathcal R_t=\frac4{22500t}.
 \label{eq:m14-counterexample-risk}
\end{equation}
For example, this follows also from $T=0$, $q=\epsilon^{-2}$,
$\delta_D=0$, and $\|e\|^2=4/22500$ in Proposition M14.1.
At $\kappa=10$,
\begin{equation}
 \lim_{\epsilon\downarrow0}V=\frac{72}{205},\qquad
 \lim_{\epsilon\downarrow0}V_{\rm pred}=0.
 \label{eq:m14-counterexample-value}
\end{equation}
The direction standard deviation here is $1/100$ radian per axis.
More generally, with fixed shading $s$ and tangent response $u$,
$\mathcal R_1/j_1\to w u^2s_D^2$ as $s_I\downarrow0$.
This can be arbitrarily large with increasing positive $w$, showing why
small degrees and an exact gauge alone cannot certify small relative error.
Conversely, when \emph{all} channel deviations tend to zero,
$J_a(t)\to a^\top D^{-1}a$ and $j_t\to r$; the spatial residual
$\delta_D$ still need not vanish.

\subsection{Numerical validation and scope of the evidence}
\label{subsec:m14-evidence}

The mathematical results above are proved independently of this experiment.
The new research runner reconstructs the unchanged nominal scenes using the
original object-specific seeds, 48 active lights, 1200 sampled pixels, and
the corrected noise-fit convention.  It reads the unrounded risk endpoints
and value fields of all 44 frozen goal-oriented cells.  Exact observed risks
are recomputed independently through the published low-rank function;
certificate construction occurs before those risk evaluations.
The implementation uses nuisance-space spectral bins, not $F(t)^{-1}$.
It also evaluates \emph{all} 45 existing C9 grid points on all 11 objects,
including all heterogeneous/coupled combinations and all 88 direction-only
controls.  The combined 539 records include 11 repeated nominal-level-0.5/C9
anchor conditions: there are 528 distinct object--parameter conditions, not
539 independent samples.  Tags also overlap, so tag-specific counts must
not be added together.
No grid points, bin widths, or the requested tightness threshold of 10 are
chosen by observed success.  A one-object exploratory calculation preceded
implementation; this extension is not represented as a prospectively
preregistered study.

The ratios in the following table are certificate upper bound divided by
observed absolute approximation error; medians, 95th percentiles, and maxima
are across the indicated endpoints or cells, without denominator flooring.
\begin{center}
\begin{tabular}{lrrrr}
\hline
Set and quantity & Count & Median & 95th percentile & Maximum\\
\hline
Frozen nominal, risk endpoints & 88 & 1.353238 & 1.546959 & 1.652930\\
Frozen nominal, value & 44 & 1.298290 & 1.435463 & 1.483561\\
Full C9 family, risk endpoints & 990 & 1.317729 & 1.535014 & 1.724254\\
Full C9 family, value & 495 & 1.317539 & 2.480009 & 4.146191\\
\hline
\end{tabular}
\end{center}
All these tested ratios are at most 10, with no bound violations.
This is not solely a high-corruption observation: at nominal level $0.1$,
all 22 endpoint ratios and all 11 value ratios satisfy the threshold;
their maxima are 1.454136 and 1.305742, respectively.  The maximum observed
value error there is 0.038539 and its largest certified upper bound is
0.045102.  Endpoint relative errors themselves need not be tiny: a tight
certificate is not an assertion that the approximation is uniformly accurate.

The largest direction-only value error is at
\texttt{obj\_19\_cylinder}, $s_I=10^{-6}$ and direction standard deviation
$0.1$ degree.  Its observed $V-V_{\rm pred}=0.4641367354$ lies in the
\emph{design-only signed interval}
$[0.2566364348,0.9046182854]$, a bound/error ratio of 1.949034.
The lower endpoint excludes 0.05 accuracy; this is stronger than merely
failing to certify it.  On the full direction-only grid, 42 cells have
observed error above 0.05 and the interval conclusively excludes 0.05
accuracy for 32 of them; the other ten remain inconclusive rather than
falsely certified.  These controls use C9's finite positive-definite
$s_I=10^{-6}$ convention, not a pseudoinverse at $s_I=0$.

This experiment evaluates exact-arithmetic certificates in double precision;
it is not a directed-rounding computer-assisted proof.  The largest
identity-to-independent-risk relative discrepancy is
$5.72\mathbin{\times}10^{-9}$; the frozen nominal endpoint replay discrepancy
is at most $4.82\mathbin{\times}10^{-10}$.  Twenty-five negative Gram
eigenvalues attributed to roundoff are recorded, with largest magnitude
$2.83\mathbin{\times}10^{-11}$ (largest relative magnitude
$7.41\mathbin{\times}10^{-18}$); larger PSD violations raise an error rather
than add a ridge.  The predeclared floating-point check tolerance is
$2\mathbin{\times}10^{-11}+2\mathbin{\times}10^{-7}|J|$ and does not alter the
certificate or tightness ratio.  The implementation reports this limitation
rather than upgrading numerical verification into an exact proof.
The prior-PSD guard propagates the parent-Gram tolerance, rather than scaling
relative to a projection Gram that may be exactly zero: if the accepted
$\widehat T\succeq-\tau\|M\|_F I$, then
\[
 Q^\top C^\top\widehat T C Q\succeq
 -\tau\|M\|_F\|C\|_2^2 I
 =-\tau\|M\|_F\lambda_{\max}(\Sigma_0)I.
\]
This follows by evaluating the quadratic form and using $\|Qx\|=\|x\|$.
Both the actual projection-Gram scale and this propagated tolerance are
reported.  It avoids falsely rejecting valid $B=AW$ models with $T=0$ after
floating-point Gram cancellation; it does not add a prior or Fisher ridge.

The companion evidence index identifies the machine-readable artifact,
its nominal and family summaries, the low-level cases, the worst direction-only
case, and the complete per-cell inputs and certificates.
All 75 pre-existing result files retain their before/after SHA-256 hashes.
The family-ranking diagnostics
$S_{\rm pred}$ and $S_{\rm real}$ concern different reconstruction endpoints
and are not redefined as remainder metrics here.  M14 establishes a
conditional fixed-linearization risk certificate, not a nonlinear recovery,
ranking-transfer, or geometry-free parameter-ratio theorem.
